\chapter{Effect estimation, heterogeneous actions and threshold comparisons}
\label{exec17:s4}

A servicing team may want to know the value of offering assistance to every
eligible account, the value among accounts already receiving assistance, or
the value in the population where receipt remains uncertain. Those questions
need not have the same answer. Nor does learning a large conditional benefit
identify the best use of a scarce specialist team. A threshold experiment can
support a local comparison while leaving the portfolio decision unresolved.
The proposal therefore needs a vocabulary for comparing estimators through
their population, uncertainty and decision consequences.

This chapter develops three connections: assignment-dependent populations and
orthogonal estimation; joint heterogeneous comparisons and resource allocation;
and threshold identification versus finite approximation. The numerical laws
are stipulated, reproducible examples in USD per eligible account over a fixed
thirty-day horizon. They establish algebra and decision consequences under
those laws. They are neither bank observations nor new executions of the
external software discussed below.~\citep{exec17_s4_math}

The later comparisons follow these distinctions into individual software
paths: the population actually averaged, the provenance of nuisance predictions,
and the construction of the final uncertainty. Repeated examples use the same
stipulated customers where that makes a difference between implementations
visible. The alternatives remain separate when their priors, sampling units,
missingness rules or numerical conventions answer different questions.

\section{The population can depend on the assignment mechanism}
\label{exec17:s4-target}

Let $O=(X,A,Y)$ denote one independent account observation, where $X$ is
information available before a binary action $A$, and $Y$ is the subsequent
monetary outcome. Assume consistency, no interference, conditional
exchangeability of the two potential outcomes given $X$, and positive
probabilities for both actions on the population being studied. These are
scientific restrictions; fitting a flexible regression does not verify them.
Define
\[
 e(x)=P(A=1\mid X=x),\qquad m_a(x)=E(Y\mid A=a,X=x),\qquad
 \tau(x)=m_1(x)-m_0(x).
\]
Throughout the following derivations, the moments exist and the denominators
are nonzero. Uniform overlap and suitable higher moments will be imposed when
an asymptotic variance or remainder bound needs them.

A useful family of targets is
\begin{equation}
 \theta_h(P)=\frac{E\{h(e(X))\tau(X)\}}{H},\qquad
 H=E\{h(e(X))\}>0.
 \label{exec17:s4-weight-target}
\end{equation}
The cases $h=1$, $h(e)=e$, $h(e)=1-e$ and $h(e)=e(1-e)$ give the population,
treated, untreated and overlap-weighted effects. The treated case follows
because $E\{e(X)f(X)\}=P(A=1)E\{f(X)\mid A=1\}$. Overlap weighting emphasizes
accounts with uncertain receipt; it does not estimate the value of treating
every account. The bank must decide whether that emphasis describes an
actionable population or only a useful scientific comparison.

There is a further distinction within the word ``weighted.'' A weight
$w(x)$ supplied independently and held fixed defines
$E\{w(X)\tau(X)\}/Ew(X)$. A weight computed from the unknown assignment law
defines a functional that changes when that law changes. Treating an estimated
overlap weight as fixed in a variance calculation can discard a component of
uncertainty. A randomized experiment with a known, fixed assignment design is
yet another model: its allowed perturbations do not include arbitrary changes
to that assignment law. The relevant derivative must match the model being
claimed.~\citep{exec17_s4_math,dml2018}

To see the missing component, consider a smooth parametric path through the
observable law, with mean-zero score $s(O)$. A dot denotes a path derivative.
Conditional differentiation gives
\begin{align*}
 \dot e(x)&=E\{(A-e(x))s(O)\mid X=x\},\\
 \dot m_a(x)&=E\{(Y-m_a(x))s(O)\mid A=a,X=x\}.
\end{align*}
The marginal distribution of $X$ also moves. Applying the quotient rule to
Equation~\eqref{exec17:s4-weight-target} separates that movement from changes
in $m_a$ and $e$. For differentiable $h$, a mean-zero influence function in
the unrestricted observational model is
\begin{align}
 \phi_h(O)={}&\frac{h(e)}{H}
 \left\{\tau-\theta_h+\frac{A}{e}(Y-m_1)
                 -\frac{1-A}{1-e}(Y-m_0)\right\}\notag\\
 &+\frac{h'(e)}{H}(A-e)(\tau-\theta_h).
 \label{exec17:s4-weight-if}
\end{align}
Functions of $X$ in this expression are evaluated at the observed $X$.
Taking its covariance with $s$ recovers the three derivative contributions.
The last line is the assignment-law contribution. It disappears for fixed
external weights, but generally survives for treated and overlap targets.
This is a project derivation under the stated model, rather than an inference
that any package implements every member of this family.~\citep{exec17_s4_math}

The extra term is orthogonal to the first line. Its conditional mean given
$X$ is zero; its covariance with each outcome residual is zero because
$E(Y-m_a\mid A=a,X)=0$. Consequently its addition increases the variance by
\begin{equation}
 \frac{E\{h'(e)^2e(1-e)(\tau-\theta_h)^2\}}{H^2}.
 \label{exec17:s4-extra-variance}
\end{equation}
This comparison concerns the two displayed random variables in this
observational model. It does not claim that every estimator using a different
model must have that variance, or that weights learned on an independent
sample have the same inferential target.

For $h(e)=e$, write $p=P(A=1)$. Combining terms in
Equation~\eqref{exec17:s4-weight-if} yields
\begin{equation}
 \phi_{\rm ATT}=
 \frac{A}{p}(Y-m_0-\theta_{\rm ATT})
 -\frac{(1-A)e}{p(1-e)}(Y-m_0).
 \label{exec17:s4-att-if}
\end{equation}
The treated-arm outcome regression cancels. This agrees with the distinct
ATE and ATTE scores developed by \citet{dml2018}; it explains why an adapter
must preserve the target-specific score and nuisance requirements. Simply
relabeling an ATE result as ATT, or retaining ATE nuisance diagnostics as if
they were necessary and sufficient for ATT, changes the mathematical claim.

Hard trimming requires additional care. With
$h(e)=1\{\ell\leq e\leq u\}$, the boundary can move under a path, and the
ordinary smooth derivative above is not automatically available. A declared
fixed subpopulation, a smooth weight, and a population selected by estimated
propensity thresholds are three different specifications. Mass near the
thresholds, whether selection uses an independent sample, and whether the
estimand conditions on the selected set matter. Removing unstable accounts
from a calculation does not by itself justify an interval for the original
population.

\section{An overlap score and its exact remainder}
\label{exec17:s4-remainder}

Orthogonality is a local derivative statement. It is useful precisely because
it can suppress first-order nuisance error, but that statement must not be
silently strengthened into consistency under arbitrary misspecification.
Consider the overlap case $h(e)=e(1-e)$. Let $q(x)$ be a fitted propensity and
$\widehat m_a(x)$ fitted arm means, evaluated on observations excluded from
their training folds. Define
\begin{align}
 D(O;q)&=A(1-q)^2+(1-A)q^2,\notag\\
 N(O;q,\widehat m)&=D(O;q)(\widehat m_1-\widehat m_0)
      +A(1-q)(Y-\widehat m_1)
      -(1-A)q(Y-\widehat m_0),\notag\\
 \widehat\theta&=\frac{\mathbb P_nN}{\mathbb P_nD}.
 \label{exec17:s4-overlap-ratio}
\end{align}
At the truth, $E D=H$. Also
$D=e(1-e)+(1-2e)(A-e)$ for binary $A$, so the centered ratio score
$(N-\theta_hD)/H$ equals Equation~\eqref{exec17:s4-weight-if}.
The assignment derivative has been incorporated into both numerator and
denominator; it is not an optional standard-error correction applied after
estimating a different target.

The bias can be computed without an asymptotic approximation. Put
$\delta=q-e$ and $\delta_a=\widehat m_a-m_a$. Conditional expectation gives
\begin{align}
 E(D\mid X)&=e(1-e)+\delta^2,\notag\\
 E(N-\theta_hD)&=
 E\left[\delta^2(\tau-\theta_h)
       +\delta\{q\delta_1+(1-q)\delta_0\}\right].
 \label{exec17:s4-overlap-remainder}
\end{align}
For example, the coefficient on $\delta_1$ before simplification is
$E(D\mid X)-e(1-q)=q(q-e)$; the coefficient on $\delta_0$ is
$-E(D\mid X)+(1-e)q=(1-q)(q-e)$. These identities explain both the
second-order propensity term and the product term. The population fitted
ratio differs from $\theta_h$ by the second line divided by
$H+E\delta^2$.~\citep{exec17_s4_math}

If $q=e$, the remainder vanishes even with inaccurate arm means. If the arm
means are correct but $q$ is arbitrary, the first term generally remains.
Thus this corrected overlap construction is locally orthogonal while lacking
global double robustness in that second direction. For a fixed external
weight and an ATE-style augmented score, the usual product remainder is a
different calculation. The word ``augmented'' cannot resolve the distinction.

Under bounded $\tau-\theta_h$ and fitted probabilities in $[0,1]$, a bound
on the absolute numerator error is
\[
 \|\tau-\theta_h\|_\infty\|q-e\|_2^2+
 \|q-e\|_2\{\|\widehat m_1-m_1\|_2+
                          \|\widehat m_0-m_0\|_2\}.
\]
Sufficient conditions for this bound to be $o_p(n^{-1/2})$ include a
squared-propensity-error condition and a propensity/outcome product-rate
condition, with a denominator bounded away from zero. An asymptotic linear
expansion also needs control of the fitted score in $L_2$, moments, and the
fold evaluation terms. Cross-fitting addresses training/evaluation dependence;
it neither proves those rates nor supplies exchangeability. The conditions in
the DML paper similarly accompany its own specified scores and statistical
model. They are not a theorem for an arbitrary callable score.~\citep{dml2018,doubleml2022}

In a bank proposal these conditions create a useful design decision. An
overlap target can avoid assigning large inverse weights to nearly deterministic
receipt, but the target itself uses the receipt mechanism. An independently
specified segment can have a simpler fixed-weight target while retaining
difficult support. Neither option makes the scientific population choice
unnecessary. The appropriate comparison holds the target fixed first and
then studies estimation error, rather than treating every numerically stable
answer as an improvement.

\section{Clipping, population change and apparent precision}
\label{exec17:s4-clipping}

Consider three pre-action strata with the following exact law:
\[
\begin{array}{c|rrrr}
 X& P(X)&e(X)&m_0(X)\ ({\rm USD})&\tau(X)\ ({\rm USD})\\ \hline
 1&1/4&1/10&30&-4\\
 2&1/2&1/2&50&4\\
 3&1/4&9/10&70&20
\end{array}
\]
Within each stratum the action is generated independently of potential
outcomes. The ATE is $6$, ATT is $54/5$, and ATC is $6/5$ USD. Those basic
target differences motivate the example; the new issue is what fitted
weight changes do to its score and uncertainty. The true overlap mass is
$H=17/100$, giving $\theta_h=86/17$ USD. True-propensity trimming to
$[1/5,4/5]$ retains only the middle stratum and gives $4$ USD.
\citep{exec17_s4_math}

Now clip the two extreme propensities to obtain
$q=(1/5,1/2,4/5)$. The unnormalized plain inverse-weighted contrast
\[
 E\left\{\frac{AY}{q}-\frac{(1-A)Y}{1-q}\right\}
\]
equals $107/8$ USD, despite an ATE of $6$. Its coefficients on baseline
outcome levels no longer cancel in each stratum. This is a population
calculation, so more observations do not remove the discrepancy. An ATE
augmented score with correct arm means still gives $6$; its residual terms
have conditional mean zero. That result does not apply to every weighted
target simply because the same fitted $q$ is used.

If one instead freezes the weights $q(1-q)$ and averages the true effects,
the target is $228/41$ USD. The corrected overlap ratio in
Equation~\eqref{exec17:s4-overlap-ratio}, still using correct arm means,
gives $36/7$ USD. Its bias relative to the true overlap effect is $10/119$.
All four quantities---ATE, true overlap, frozen fitted-weight average and
corrected fitted ratio---can be estimated precisely while answering different
questions. With the stipulated arm-mean errors
$\delta_1=(1,-2,3)$ and $\delta_0=(-1,1,2)$ USD, the corrected ratio's
unnormalized error is $-239/3400$ USD and its denominator is $7/40$;
the normalized error is $-239/595$ USD. The exact decomposition therefore
checks the sign and size of a genuine product remainder, rather than merely
showing that a program terminates.~\citep{exec17_s4_math}

The familiar effective-sample-size diagnostic can also look better after
clipping. For the treated weight $W=A/q(X)$, its population analogue is
\[
 \frac{\operatorname{ESS}}{n}\longrightarrow
       \frac{\{E(W)\}^2}{E(W^2)}.
\]
Under the true propensity it is $9/34$; under the clipped propensity it is
$841/2024$. The larger second fraction measures concentration of those
weights. It neither proves target preservation nor yields an outcome-specific
standard error: outcome noise, effect heterogeneity and fitted-score error
still matter. A self-normalized estimator would be a further algorithm with
its own population ratio, not a justification for the unnormalized contrast
used above.

The assignment derivative can be checked directly in this same law. Perturb
the three propensities by $t(1/20,-1/10,1/20)$, holding stratum masses and
outcome laws fixed. The ATT derivative is $27/50$ USD per unit $t$ and the
overlap-effect derivative is $-24/17$. The covariance of the first line of
Equation~\eqref{exec17:s4-weight-if} alone with the assignment score is zero;
the complete influence function gives exactly these derivatives. Mean zero
of a candidate influence function is therefore insufficient. The derivative
test detects what a simple centering test misses.~\citep{exec17_s4_math}

\section{Selected software paths embody different estimators}
\label{exec17:s4-paths}

The inspected DoubleML paper describes linear score architecture and
model-specific nuisance estimation, with extension routes for callable scores
and sample splitting. The inherited local execution used the selected IRM
ATE path in version 0.11.4, released 10 August 2026, with externally supplied
nuisances and folds, no IPW normalization and truncation at 0.01. That is
evidence for its recorded operation only. The inspected source adjusts fitted
propensities before score evaluation. Changing truncation, target or callable
score would require a new comparison, including the remainder appropriate to
that functional.~\citep{doubleml2022,runtime_doubleml_scope,exec17_s4_paths}

The pinned R \texttt{grf} implementation takes a different route for its
binary-treatment overlap target. Its average-effect source fits a regression
of outcome residuals on treatment residuals, with an intercept, and uses a
heteroskedastic or cluster sandwich calculation. It does not simply run the
all-population AIPW average with another averaging weight. At true nuisance
functions, write $\mu(X)=E(Y\mid X)$ and
$Y-\mu=(A-e)\tau+\varepsilon$, with $E(\varepsilon\mid A,X)=0$.
The residual-regression slope is
\[
 \frac{E\{(A-e)(Y-\mu)\}}{E\{(A-e)^2\}}
       =\frac{E\{e(1-e)\tau\}}{H}.
\]
Because both residuals have mean zero at truth, an intercept does not change
this population slope. Its first-order residual-regression influence function
at truth is
\[
 \frac{(A-e)^2(\tau-\theta_h)+(A-e)\varepsilon}{H},
\]
which equals the full overlap influence function in
Equation~\eqref{exec17:s4-weight-if} for binary $A$.
This algebra connects the selected implementation to the assignment-dependent
target.~\citep{exec13_grf_rpath,exec17_s4_math}

Equality at truth does not make fitted algorithms identical. In the clipped
three-stratum example, retain the true marginal outcome mean $\mu$ but use
$q$ for treatment residuals. Here $E(A-q)=0$,
$E(q-e)^2=1/200$, and the residual-regression slope is
\[
 \frac{E\{e(1-e)\tau\}}{H+E(q-e)^2}=\frac{172}{35}.
\]
It differs both from the true overlap effect $86/17$ and from the corrected
ratio $36/7$ using correct arm means. These nuisance specifications are
deliberately different: true $\mu$ is not the same fitted object as the two
true arm means combined with an incorrect $q$. The example shows why a
common target and common limiting influence function do not establish
finite fitted equivalence.~\citep{exec17_s4_math}

At revision \texttt{5bee99b5}, the selected \texttt{grf} source's overlap branch
does not use the method choice after entering that branch. Earlier guards
still reject particular TMLE calls with clustering or sample weights. Thus
``method ignored'' must be qualified by control flow. Other inspected forest
types also have restrictions on target and method combinations. An adapter
should validate the exact combination, preserve package diagnostics and
record the actual route used. The generalized-forest paper's pointwise
normality result is conditional on its assumptions and honest, regular,
subsampled forest construction; it does not establish inference for every
selected policy or every fitted forest.~\citep{exec13_grf_rpath,exec13_grf_native,grf2018}

These comparisons favor wrapping selected public interfaces before porting
their internals. Reproducing an R or Python score in a common C++ core could
eventually reduce interprocess overhead, but would create an additional
statistical implementation to maintain. First establish the mathematical
target, nuisance convention, supported uncertainty unit and failure behavior
on matched examples. Reject unsupported target/method combinations at that
boundary. No current source inspection establishes complete dependency rights
or the performance of a new wrapper.~\citep{exec17_s4_design}

\section{Shared controls make heterogeneous comparisons joint}
\label{exec17:s4-joint}

Suppose the bank can choose no contact, a message or specialist assistance,
indexed by $A=0,1,2$. Let $e_j(X)$ and $m_j(X)$ denote action probabilities
and arm means. Under the corresponding multiarm consistency, exchangeability
and support assumptions, define every contrast on the same eligible-account
law:
\[
 \theta_j=E\{m_j(X)-m_0(X)\},\qquad j=1,2.
\]
Its population influence function is
\[
 \phi_j=\tau_j(X)-\theta_j+
       \frac{1\{A=j\}}{e_j}(Y-m_j)
       -\frac{1\{A=0\}}{e_0}(Y-m_0).
\]
For distinct noncontrol actions $j,k$, mutually exclusive arm indicators
eliminate their outcome-residual cross-products, but the shared control
residual appears twice. Therefore
\begin{equation}
 \operatorname{Cov}(\phi_j,\phi_k)=
 \operatorname{Cov}\{\tau_j(X),\tau_k(X)\}
       +E\{\sigma_0^2(X)/e_0(X)\},
 \label{exec17:s4-shared-control}
\end{equation}
where $\sigma_0^2(X)=\operatorname{Var}(Y\mid A=0,X)$.
Conditional mean-zero residuals remove the remaining cross terms. The first
term can have either sign; a positive shared-control term does not imply
positive total covariance in every heterogeneous population.

Take constant effects, assignment probabilities $(1/2,1/4,1/4)$, and
conditional outcome-noise variance $4$ USD$^2$ in all arms. The influence
covariance matrix and variance of the difference are
\[
 \Sigma=\begin{pmatrix}24&8\\8&24\end{pmatrix},\qquad
 \operatorname{Var}(\phi_2-\phi_1)=32.
\]
Adding marginal variances as though the contrasts were independent would give
$48$. In this example the shared control cancels in the difference. For an
iid mean, the corresponding asymptotic covariance is $\Sigma/n$; these are
not empirical standard errors and are not formulas for clustered bank
histories. Joint confidence regions or multiple-comparison decisions require
an explicitly chosen simultaneous procedure. Two marginal intervals alone
do not specify one.~\citep{exec17_s4_math}

This issue is concrete in the pinned CausalML T-learner path. Its point mean
uses the full prediction covariate matrix, while the inspected analytic
variance uses the observations in the current treatment/control pair. Under
unequal assignment or covariate composition, those sets can represent
different populations. The finding is about the selected revision and path,
not a claim that every CausalML estimator behaves that way. A proposed
multiarm wrapper should standardize every action prediction to one explicit
evaluation population, then construct the joint score and covariance there.
It should keep package-specific uncertainty as a separately named result
until equivalence is established.~\citep{exec14_causalml_primary,exec17_s4_paths}

The same source offers a row bootstrap. Rows from repeated accounts or a
shared calendar shock are not independent sampling units simply because they
have distinct row indices. Whether the comparator uses independent accounts,
households, branches or time blocks belongs in the statistical specification.
A bootstrap setting cannot repair an undeclared unit or create joint support
for an action rarely observed in a subgroup.

\section{A resource constraint changes the optimal action}
\label{exec17:s4-allocation}

Let $g(x)=\tau(x)-c(x)$ be the incremental net value of one action, where
$c(x)$ includes the action's incremental cost in the same USD units and
horizon. Let $r(x)>0$ be specialist resource units and let $d(x)\in[0,1]$
be the fraction of that stratum assigned to the action. With expected
resource budget $B$ per eligible account, consider
\begin{equation}
 \max_d E\{g(X)d(X)\}\quad\text{subject to}\quad
 E\{r(X)d(X)\}\leq B,\quad 0\leq d\leq1.
 \label{exec17:s4-resource-problem}
\end{equation}
Fractional decisions describe randomization or divisible allocation in this
population problem. They are not a guarantee that a finite day's integer
roster fits a hard capacity constraint.

For a multiplier $\lambda\geq0$, every feasible policy obeys the dual bound
\begin{equation}
 E(gd)\leq \lambda B+E\{(g-\lambda r)_+\}.
 \label{exec17:s4-resource-dual}
\end{equation}
Indeed, $gd=(g-\lambda r)d+\lambda rd$ and
$(g-\lambda r)d\leq(g-\lambda r)_+$. Equality is attained if the policy
uses the action wherever $g>\lambda r$, avoids it wherever $g<\lambda r$,
and allocates any equality stratum to meet the budget when the multiplier is
positive. This is a checkable optimality certificate. It explains why ranking
by effect, or even by net gain without resource units, can be wrong.

Consider four strata:
\[
\begin{array}{c|rrrr}
 \text{stratum}&1&2&3&4\\ \hline
 P(X)&2/5&3/10&1/5&1/10\\
 \text{gross effect, USD}&4&9&15&25\\
 \text{incremental cost, USD}&1&3&7&14\\
 g(X),\ \text{USD}&3&6&8&11\\
 r(X),\ \text{units}&1&1&2&4
\end{array}
\]
Set $B=1/2$ unit per eligible account. The policy
$d=(0,1,1/2,0)$ uses exactly $1/2$ unit and produces $13/5$ USD of net value.
At $\lambda=4$ USD per resource unit, the adjusted gains
$g-\lambda r$ are $(-1,2,0,-5)$; the dual bound is
$4(1/2)+(3/10)2=13/5$. The feasible allocation therefore attains the global
bound. Sorting by gross effect or by unadjusted net gain fills stratum 4
first and then a quarter of stratum 3, giving
$d=(0,0,1/4,1)$ and $3/2$ USD. The loss is $11/10$ USD per eligible
account, even though the ranked effects are known exactly.~\citep{exec17_s4_math}

This example deliberately separates scientific estimation from operational
optimization. The cost schedule is stipulated, specialist time is a separate
constraint, and the value is incremental relative to no action. If monetary
cost already prices the same scarce time, the sponsor must explain whether
the hard constraint remains necessary and avoid counting the same resource
twice. Risk limits, fairness conditions and action eligibility can further
restrict the feasible set. For indivisible actions on a finite roster, the
appropriate optimization can become a knapsack or assignment problem; the
fractional certificate does not assert that a greedy integer solution is
optimal.

For multiple actions, replace $d$ by $d_j(x)\geq0$ with
$\sum_jd_j(x)\leq1$, action-specific gains and resources. At a fixed
multiplier, each account receives the action with the largest positive
$g_j-\lambda r_j$, subject to any additional declared constraints. Shared
controls and common target populations remain essential because the optimizer
compares these action values directly. Exporting unrelated pairwise targets
to an optimizer can produce a feasible but scientifically meaningless plan.

\section{From effect error to decision loss}
\label{exec17:s4-regret}

Prediction error and policy error need not move together. For the same fixed
feasible policy class $\mathcal D$, let $d^*$ maximize $E(gd)$ and
$\widehat d$ maximize $E(\widehat g d)$. Suppose the pointwise bound
$|\widehat g-g|\leq\epsilon(X)$ is valid. Adding and subtracting the fitted
objective gives
\begin{align}
 E\{g(d^*-\widehat d)\}
 &=E\{(g-\widehat g)(d^*-\widehat d)\}
       +E\{\widehat g(d^*-\widehat d)\}\notag\\
 &\leq E\{\epsilon(X)|d^*-\widehat d|\}.
 \label{exec17:s4-regret-bound}
\end{align}
The last fitted-objective difference is nonpositive by optimality. If the
optimizer is within $\eta$ USD of its optimum, add $\eta$ to the bound.
This is a deterministic comparison conditional on the two gain functions,
the common population objective and the same feasible class. It is not a
confidence bound supplied by an arbitrary learner. If the fitted optimizer
uses an empirical objective, a population-to-empirical discrepancy must also
be controlled. If costs or capacity are estimated, their errors can change
the feasible set, so this derivation no longer covers them by itself.

With $0\leq d\leq1$, a uniform gain-error bound $\epsilon_0$ gives the coarse
bound $\epsilon_0$ per account. The sharper expression charges error only
where the two policies disagree. This explains why errors near a resource
threshold can be more consequential than equally large errors in strata
where the same action remains optimal. A low global squared prediction loss
does not alone demonstrate low constrained decision loss. The finite
reference enumerates 81 signed-error configurations over the binary feasible
policies and confirms the inequality, while the algebra establishes its stated
general domain.~\citep{exec17_s4_math}

The observational policy-learning analysis of \citet{policy2021} supplies a
more substantive statistical route: construct doubly robust scores and learn
within a controlled-complexity policy class under overlap, noise and nuisance
rate conditions. Its regret theorem is tied to the admitted binary-decision
problem and assumptions. It does not automatically cover every resource
constraint, multiarm action, dependent history or optimizer introduced here.
Likewise, a pointwise conditional-effect interval from a forest is not an
interval for the value of a policy selected using the same sample.

A proposal comparison should therefore split learning and evaluation at the
appropriate independent unit, declare the reference population and feasible
class, and evaluate the frozen selected policy with a justified value score.
It should report infeasibility and inability to estimate value as outcomes.
An exact reference can verify the optimizer; a prospective generated-outcome
study can investigate selection and finite-sample inference. Neither supplies
the true net effect of an untested bank action.

\section{Threshold estimates include an approximation problem}
\label{exec17:s4-curvature}

Now let $S$ be a pre-action eligibility score, centered at a cutoff of zero.
A sharp rule assigns the action when $S\geq0$. Continuity of the relevant
potential-outcome conditional means at zero, suitable support approaching
the cutoff and a well-defined assignment rule can identify a local cutoff
effect. Sorting, manipulation, concurrent threshold rules and delivery
versions can invalidate that interpretation. The following calculation
assumes the scientific cutoff question has already been justified and asks
what a finite-bandwidth approximation estimates.

On the right side write $S=hu$, with $u\in[0,1]$, and take uniform local
design density and uniform weights for an exact population local-linear
projection. Suppose the side's conditional mean is
$b_0+b_1S+\kappa S^2$. The normal equations are
\[
 G\begin{pmatrix}\alpha\\\beta\end{pmatrix}
 =\int_0^1\begin{pmatrix}1\\u\end{pmatrix}
       (b_0+b_1hu+\kappa h^2u^2)\,du,
 \qquad
 G=\begin{pmatrix}1&1/2\\1/2&1/3\end{pmatrix}.
\]
Since $G^{-1}=\left(\begin{smallmatrix}4&-6\\-6&12\end{smallmatrix}\right)$,
the intercept equals
\begin{equation}
 \alpha=\int_0^1(4-6u)m(hu)\,du
        =b_0-\frac{\kappa h^2}{6}.
 \label{exec17:s4-local-bias}
\end{equation}
The equivalent kernel integrates to one and has zero first moment, but its
second moment is $-1/6$. Its negative weights beyond $u=2/3$ are part of
boundary trend correction. They are not positive population weights.
Here $\kappa$ is the quadratic coefficient; the second derivative is
$2\kappa$. Confusing the two would double the bias correction.
\citep{exec17_s4_math}

Using $S=-hu$ on the left changes the linear term's sign and leaves the
quadratic intercept bias unchanged. With equal bandwidths, the difference
between projected intercepts has bias
$-(\kappa_+-\kappa_-)h^2/6$. Unequal bandwidths replace this by
$-(\kappa_+h_+^2-\kappa_-h_-^2)/6$. These are exact projection identities
for the stipulated quadratic and uniform-density law. A general density,
another kernel or a nonquadratic regression requires its own expansion and
remainder. A finite random-sample local regression is not identical to the
population projection; its expected ratio of random matrices is another
object.

Reducing bandwidth can reduce approximation bias while removing nearby
observations and destabilizing the local design. Estimated curvature,
bandwidth selection and robust bias correction add their own uncertainty.
An interval using a variance formula for an uncorrected intercept cannot be
declared valid merely because a point bias adjustment was subtracted. The
inspected \texttt{rdrobust} help documents local-polynomial, bias-corrected
and variance choices; the formula above is a transparent reference for a
particular population calculation, not a substitute for its full inference
theory.~\citep{exec13_ivrd_rdrobust}

\section{A fuzzy ratio can amplify approximation error}
\label{exec17:s4-fuzzy}

With incomplete receipt, let $\Delta_Y$ and $\Delta_A$ denote the true
outcome and receipt discontinuities. Under the separately required exclusion,
monotonicity and local continuity restrictions, their ratio
$\theta=\Delta_Y/\Delta_A$ has a local complier interpretation when the
receipt jump is nonzero. Let $b_Y,b_A$ be the biases in the population
projected discontinuities. Exact algebra gives
\begin{equation}
 \frac{\Delta_Y+b_Y}{\Delta_A+b_A}-\theta
       =\frac{b_Y-\theta b_A}{\Delta_A+b_A}.
 \label{exec17:s4-fuzzy-bias}
\end{equation}
Even before sampling uncertainty, an apparently small bias in the receipt
jump can change the denominator's magnitude or sign. Correcting only the
outcome projection does not generally correct the ratio.

Set the true local effect to $8$ USD, the outcome curvature-coefficient
difference to $12$ USD per score-unit squared and the receipt difference to
$1/2$ per score-unit squared. Then $b_Y=-2h^2$ USD and $b_A=-h^2/12$.
The exact projected ratios are
\[
\begin{array}{c|cc|c}
 \Delta_A&h&\Delta_A+b_A&
    (\Delta_Y+b_Y)/(\Delta_A+b_A)\ ({\rm USD})\\ \hline
 1/2&1/2&23/48&168/23\\
 1/2&1/4&95/192&744/95\\
 1/50&1/2&-1/1200&408\\
 1/50&1/4&71/4800&168/71
\end{array}
\]
All rows have the same true ratio of $8$. At the weak receipt jump, the
first projected ratio is $408$ USD, although no sampling noise has been
introduced. At the narrower bandwidth it is approximately $2.37$ USD.
These values expose approximation interacting with a weak denominator.
They are not expectations of an upstream ratio estimator and do not show
that an actual package returns these values.~\citep{exec17_s4_math}

Knowing the exact curvature in this artificial law would permit exact
population correction. In data, estimating curvature creates another random
quantity, and a weak first stage remains a sampling problem after correction.
The product should distinguish an approximation certificate, an interval for
the ratio, and a scientific assessment of exclusion and local compliance.
It should preserve unbounded or disconnected uncertainty where the selected
valid procedure permits it. A finite number from a stabilized denominator is
not itself evidence for a bounded causal effect.

\section{Finite contrasts, score support and portfolio transport}
\label{exec17:s4-support}

A second approximation arises when software evaluates predictions at nearby
scores. In the selected pinned CausalPy regression-discontinuity path, the
contrast uses design rows for treatment above $+\epsilon$ and control below
$-\epsilon$. This is a finite contrast, not automatically the limiting
cutoff effect. For example, let the two mean functions be
\[
 m_1(s)=18+5s+12s^2,\qquad m_0(s)=10+3s.
\]
Their cutoff difference is $8$ USD, while
\[
 m_1(\epsilon)-m_0(-\epsilon)=8+8\epsilon+12\epsilon^2.
\]
At $\epsilon=1/10,1/100,1/1000$, the contrasts are $8.92$, $8.0812$ and
$8.008012$ USD. Decreasing $\epsilon$ reduces this particular deterministic
approximation. It does not create data near the cutoff or eliminate model
misspecification. A numerical tolerance should therefore be expressed in
score units and translated into a justified monetary approximation bound
when the model permits one.~\citep{exec14_causalpy_primary,exec17_s4_math}

The same CausalPy source distinguishes expected mean predictions from
posterior predictive observations. Predictive variation can include outcome
noise; posterior uncertainty in a mean contrast is a different object.
The inspected scikit-learn route's singleton draw axes are not a posterior
sample. A common result schema should state the uncertainty object and
supported backend. Admitting an array shape as evidence of a distribution
would turn a formatting convention into an unsupported scientific claim.
\citep{exec14_causalpy_primary,exec17_s4_paths}

Discrete eligibility scores can remove the limiting comparison entirely.
Suppose $S$ takes only $-1$ and $1$, action is $1\{S\geq0\}$, and observed
outcomes are respectively $10$ and $14$ USD. Two response systems agree on
every observation:
\[
 \text{system I: }Y(a)=10+4a,\qquad
 \text{system II: }Y(a)=12+2S.
\]
Both have smooth extensions as functions of the score. The cutoff treatment
effect is $4$ in the first system and $0$ in the second. Continuity alone
cannot distinguish them because the score support never approaches zero.
Fitting a line across the gap chooses an extrapolation restriction. A
mass-point diagnostic or a bandwidth adjustment can reveal or respond to
the observed support, but cannot supply the unobserved response surface.
\citep{exec17_s4_math}

Possible remedies depend on the bank's decision. A defensible local
randomization model for a specified finite window would be a different
scientific assumption. A justified functional-form or curvature restriction
could support extrapolation or bounds. A redesigned randomized encouragement
could provide additional intervention information. Without such information
or restrictions, the correct product result is an unresolved effect set,
not a universal recommendation to use a different smoother.

Even a well-identified continuous-score cutoff effect need not equal a
portfolio effect. Suppose $\tau(0)=8$ USD and stipulate the sensitivity
restriction
\[
 |\tau(s)-\tau(0)|\leq L|s|,\qquad
 L=6\ \text{USD per score unit}.
\]
For a target population with equal mass at $s=-1/2$ and $s=1/2$, integrating
the pointwise envelope gives
\begin{equation}
 5\leq E_T\{\tau(S)\}\leq11\quad\text{USD}.
 \label{exec17:s4-transport-bounds}
\end{equation}
These bounds are sharp relative to the stipulated cutoff, sensitivity
restriction and observed sharp-threshold law. To construct witnesses, take
the observed arm mean to be $10+2s+8a$ with $a=1\{s\geq0\}$. Preserve
that mean on the observed arm, set the effect to $8-6|s|$ or $8+6|s|$,
and fill in the unobserved arm mean by subtracting or adding the effect.
The two continuous response systems agree on every observed mean and at the
cutoff while attaining the two portfolio endpoints.~\citep{exec17_s4_math}

At a cost of $7$ USD per action, the portfolio net-value set is $[-2,4]$ USD.
The local cutoff net value of $1$ USD does not determine its sign. The constant
$L$ is a scientific sensitivity parameter, not an estimate certified by
the local effect. Additional justified shape restrictions, intervention data
away from the cutoff or a narrower portfolio could reduce the set. A
posterior extrapolation can also be reported, provided its model-dependent
meaning remains explicit. General policy eligibility and transport cannot be
inferred from a local ratio alone.

\section{Adoption choices and a public comparison that can be built}
\label{exec17:s4-adoption}

The practical comparison is among scientific routes as well as languages.
DoubleML offers an inspected Python score architecture; \texttt{grf} supplies
R interfaces backed by native forest components; CausalML supplies heterogeneous
prediction and selected uncertainty paths; CausalPy supplies model-based
threshold predictions; and \texttt{rdrobust} supplies specialized local
polynomial procedures. An existing native component is evidence of an
implementation route, not evidence that a newly ported common core has the
same target or interval behavior. The following decisions are prospective
and scoped to the inspected paths.~\citep{exec17_s4_paths,exec17_s4_design}

\begin{longtable}{p{.19\textwidth}p{.35\textwidth}p{.36\textwidth}}
\toprule
Candidate route&Proposed use&Condition or refusal\\ \midrule
DoubleML IRM&Wrap selected ATE/ATTE score paths; retain callable-score extension
as a separate development.&Bind folds, nuisance predictions, clipping,
normalization and target. A new overlap score needs its own derivative,
remainder and statistical study.\\
\texttt{grf}&Wrap the supported average-effect and heterogeneous-prediction
routes before considering a native port.&Preserve target/method control flow,
cluster units and residual-regression semantics. Do not equate pointwise
forest intervals with selected-policy uncertainty.\\
CausalML&Adopt selected conditional predictions as candidate nuisances or
action scores.&Build a common-population evaluator for multiarm comparisons;
verify joint covariance. Retain the selected package variance separately when
its population differs.\\
CausalPy&Use the inspected threshold construction for explicit model-based
comparisons.&Name finite $\epsilon$, mean versus predictive uncertainty and
backend. Reject an unsupported cutoff interpretation for a score gap.\\
\texttt{rdrobust}&Compare specialized sharp/fuzzy local-polynomial procedures
on their admitted designs.&Bind the exact implementation, bandwidths,
polynomial orders and variance method; distinguish mass-point handling from
scientific identification.\\
Licensed Stata host&Offer a documented host route where the bank already has
the required entitlement.&The inspected package requires Stata 16 or later.
GPL-3 package rights and commercial host rights are separate; deployment,
redistribution and cost need a matched review.\\ \bottomrule
\end{longtable}

The inspected Stata \texttt{rdrobust} help declares version 11.1.1 dated
1 October 2026, with controls for outcome and bias-polynomial orders, effect
and bias bandwidths, fuzzy receipt, mass points and variance. The inherited
repository selection used a UTC end-of-day cutoff query. Whether that exact
revision was available by the strict Shanghai client cutoff remains unresolved.
The subsequently acquired mutable Stata 19 manuals also lack certified
cutoff-date bytes. Their inspected documented procedures can inform a bounded
comparison while those identity limits remain visible; they cannot establish
an exact as-of-date commercial capability census.~\citep{exec13_ivrd_rdrobust,exec13_ivrd_stata}

The commercial comparison should therefore ask whether the licensed route
answers the same target with the same support and sampling assumptions, which
diagnostics and uncertainty procedures it documents, and what integration and
operating costs apply. The proposal's desired advantage is preservation of
these distinctions across discovery, identification, estimation and reporting.
That is a proposed capability, not a measured victory over a commercial tool.
No proprietary price, actual deployment entitlement or new vendor runtime is
established by the present source review.

A bounded public/generated prototype can test the proposed advantage without
waiting for bank data. First freeze the package revisions, environments,
complete dependency rights and target contracts. Replay the exact weighted,
multiarm, resource and threshold laws above to test population arithmetic,
score conventions and expected refusals. Use the same declared target when
comparing algorithms; deliberately changed targets form separate cases.
The current exact reference calculations perform only the mathematical
part of that work.~\citep{exec17_s4_design}

For finite-sample statistical claims, separately register a prospective study:
sampling and assignment laws, outcome noise, independent units, nuisance
learners, folds, bandwidth rules, all uncertainty procedures, seed, fixed
request counts and complete failure accounting. Public covariates can supply
realistic margins while generated potential outcomes preserve known causal
truth. Their public origin does not make a generated intervention mechanism
authentic bank evidence. Compare common-target bias and error, interval
coverage and width, availability, allocation feasibility, policy loss and full
observed computation cost. Include clipped propensities, weak receipt jumps,
discrete support and dependent-unit mismatches as declared stress regimes,
rather than silently excluding failed requests.

Successful arithmetic and statistical comparisons would justify a bounded
prototype and identify where wrapping suffices or porting merits investment.
They would leave real confounding, score manipulability, feasible bank actions,
institutional capacity, acceptable loss and actual total prices to the later
qualification decision. Where two scientific models still agree on the
observations but disagree on the effect, more software cannot resolve the
ambiguity. The product must either acquire the missing information, narrow
its assumptions or preserve the effect set through the financial decision.

The next analysis develops the shared-control, allocation and discrete-score
cases further. The source comparisons then ask whether the implementations
preserve those targets and uncertainties. This order separates what the bank
needs to estimate from what a particular numerical interface happens to return.
\label{exec17:s4-end}
